% Generated draft: requires mathematical and source review.
\documentclass[11pt]{article}
\usepackage[margin=25mm]{geometry}
\usepackage{fontspec}
\setmainfont{Latin Modern Roman}
\usepackage{amsmath,amssymb,mathtools,unicode-math}
\setmathfont{Latin Modern Math}
\usepackage{enumitem,needspace,xurl,hyperref}
\setlength{\parindent}{0pt}
\setlength{\parskip}{5pt}
\setlength{\emergencystretch}{3em}
\newcommand{\N}{\mathbb{N}}
\newcommand{\Z}{\mathbb{Z}}
\newcommand{\Q}{\mathbb{Q}}
\newcommand{\R}{\mathbb{R}}
\newcommand{\C}{\mathbb{C}}
\newcommand{\F}{\mathbb{F}}
\newcommand{\A}{\mathbb{A}}
\newcommand{\PP}{\mathbb{P}}
\newcommand{\E}{\mathbb{E}}
\newcommand{\eps}{\varepsilon}
\newcommand{\dd}{\,\mathrm{d}}
\begin{document}
\section*{UnsolvedMath ID: 30004894}
\textbf{Source problem code:} catalogue entry
\subsection{Definitions and mathematical statement}
Let \(Q\) be a quiver of affine type, and let \(Q_0\) denote its vertex set. Let \(k=\mathbb F_q\), where \(q\) is a prime power. The notation \(\delta\) conventionally denotes the minimal positive imaginary root of the corresponding affine root system; the supplied excerpt does not explicitly define it or specify the precise class of affine quivers allowed.

The supplied record calls \(C(n\delta)\) the semicuspidal algebra. Its construction, coefficient conventions, and auxiliary choices are not supplied. The excerpt also mentions a KLR algebra \(R(n\delta)\) and an ideal \(I_\mu(n\delta)\), without defining either or explicitly giving the quotient defining \(C(n\delta)\). The notation \(S_n(\mathbb Z^{Q_0})\) denotes the proposed Schur algebra in the cleaned record, but neither this construction nor the mathematical object represented by \(\mathbb Z^{Q_0}\) is defined in the supplied material.

\textbf{Conjecture 6, as represented by the supplied cleaned record.} For every quiver \(Q\) of affine type and every finite coefficient field \(k=\mathbb F_q\), the proposed identification is
\[
 C(n\delta)\simeq S_n(\mathbb Z^{Q_0})\otimes_{\mathbb Z}\mathbb F_q.
\]
Here \(n\) is the parameter occurring in the source. Its permitted range, including whether \(n=0\) is included, is not stated. The tensor product presupposes an integral algebra structure on the proposed Schur algebra, whose definition is absent. The intended assertion is an algebra isomorphism; whether it must preserve a grading or additional structure is unspecified.

\textbf{Unresolved source qualification.} The extracted original writes the right-hand algebra as \texttt{Sen (ZQ0)}, whereas the cleaned statement writes \(S_n(\mathbb Z^{Q_0})\). The excerpt separately distinguishes \(S_n(\mathbb P^1)\) from an image algebra also extracted with the notation \texttt{Sen}. Consequently, a potentially meaningful decoration on the conjectured Schur algebra cannot be resolved from this record. The displayed formulation records the cleaned target conditionally; it does not identify these potentially distinct algebras or supply their missing definitions.

The proposed geometric argument is explicitly unfinished in the supplied excerpt. Its subsequent Steinberg-variety assertion is labelled a \emph{theorem in progress}, and is not supplied as an established proof of this conjecture.
\subsection{Short English statement}
For an affine quiver, the conjecture proposes an isomorphism between its semicuspidal algebra at dimension \(n\delta\) and the reduction of a proposed integral Schur algebra modulo the characteristic of \(\mathbb F_q\). The supplied excerpt leaves essential definitions, the range of \(n\), and a potentially significant Schur-algebra decoration unresolved. This is a statement-only draft requiring human mathematical review.
\subsection{Sources}
\begin{itemize}
\item UnsolvedMath local registry, permanent ID \(30004894\), source problem code \texttt{catalogue entry}, title: \emph{Semicuspidal and Schur Algebras of Affine Quivers}. The supplied record identifies the problem as Conjecture 6, records proposed year 2021, and names no proposer.
\item Supplied source citation: \emph{Enveloping Algebras and Geometric Representation Theory} (2022), Oberwolfach Reports, pp.\ 2827--2891, DOI \texttt{10.4171/owr/2021/52}. \url{https://doi.org/10.4171/owr/2021/52}. The citation year 2022 and proposed year 2021 are preserved as supplied.
\item Supplied literature source: ``Oberwolfach Report 52, affine KLR.'' \url{https://doi.org/10.4171/owr/2021/52}.
\item Supplied literature source: ``Maksimau--Minets, KLR and Schur algebras.'' \url{https://academic.oup.com/imrn/article/2023/8/6976/6554160}.
\item Literature assessment dated exactly \texttt{2026-08-22}: the supplied curation labels the problem open, states that the report leaves the microlocal argument unfinished, and states that later related KLR and Schur-algebra work does not establish this exact conjecture. This is supplied dated curation, not an independently verified literature assessment.
\item The registry labels the statement \texttt{corrected\_verified} and reports that the official report was checked directly. These are supplied verification claims; no independent source consultation or mathematical verification is asserted here. The excerpt's reference label [4] is not bibliographically resolved in the supplied material.
\end{itemize}
\end{document}
